\paragraph{归一化 Lucas 幂 Hankel 主子式的离散曲率与 residue-lumping 的 Ramanujan 图册}
在定义 \ref{def:conclusion-lucas-fib-input} 的记号下，对 \(q\ge 0\) 约定
\[
\mathcal H_0:=1,\qquad
\mathcal H_q:=
\frac{\det(H_q)}
{\left(\prod_{k=0}^{q}\binom{q}{k}\right)\Delta^{q(q+1)/2}}
\qquad (q\ge 1),
\]
并记前向差分算子
\[
\nabla A_q:=A_q-A_{q-1},
\qquad
\nabla^2A_q:=A_q-2A_{q-1}+A_{q-2}.
\]

\begin{theorem}[归一化 Lucas 幂 Hankel 主子式的离散曲率公式]\label{thm:conclusion-lucas-normalized-hankel-discrete-curvature}
对一切 \(q\ge 1\)，有
\[
\mathcal H_q=\prod_{d=1}^{q}F_d^{\,2(q+1-d)}.
\]
进一步，对一切 \(q\ge 1\) 与 \(q\ge 2\) 分别成立
\[
\frac{\mathcal H_q}{\mathcal H_{q-1}}
=
\prod_{d=1}^{q}F_d^{\,2},
\qquad
\frac{\mathcal H_q\mathcal H_{q-2}}{\mathcal H_{q-1}^{\,2}}
=
F_q^{\,2}.
\]
\leanverified{paper\_conclusion\_lucas\_normalized\_hankel\_discrete\_curvature}
\end{theorem}
\begin{proof}
第一式只是定理 \ref{thm:conclusion-lucas-power-hankel-closed-form} 对归一化量 \(\mathcal H_q\) 的直接改写。于是
\[
\frac{\mathcal H_q}{\mathcal H_{q-1}}
=
\frac{\prod_{d=1}^{q}F_d^{\,2(q+1-d)}}
{\prod_{d=1}^{q-1}F_d^{\,2(q-d)}}
=
\prod_{d=1}^{q}F_d^{\,2},
\]
从而得到第一条商公式。再对相邻商作比，
\[
\frac{\mathcal H_q\mathcal H_{q-2}}{\mathcal H_{q-1}^{\,2}}
=
\frac{\mathcal H_q/\mathcal H_{q-1}}{\mathcal H_{q-1}/\mathcal H_{q-2}}
=
\frac{\prod_{d=1}^{q}F_d^{\,2}}{\prod_{d=1}^{q-1}F_d^{\,2}}
=
F_q^{\,2}.
\]
证毕。
\end{proof}

\begin{corollary}[Archimedean Fibonacci 对数梯的二阶恢复]\label{cor:conclusion-lucas-log-fibonacci-from-hankel-curvature}
若处于一条实分支并满足 \(F_d>0\)（对所论全部 \(d\)），定义
\[
\mathcal L_q:=\log \mathcal H_q\qquad(q\ge 0).
\]
则对一切 \(q\ge 1\) 与 \(q\ge 2\) 分别有
\[
\nabla \mathcal L_q
=
2\sum_{d=1}^{q}\log F_d,
\qquad
\nabla^2\mathcal L_q
=
2\log F_q.
\]
等价地，
\[
\log F_q
=
\frac12\bigl(\mathcal L_q-2\mathcal L_{q-1}+\mathcal L_{q-2}\bigr)
\qquad(q\ge 2).
\]
\end{corollary}
\begin{proof}
对定理 \ref{thm:conclusion-lucas-normalized-hankel-discrete-curvature} 中的商公式取对数，得到
\[
\mathcal L_q-\mathcal L_{q-1}
=
\log\!\left(\frac{\mathcal H_q}{\mathcal H_{q-1}}\right)
=
2\sum_{d=1}^{q}\log F_d.
\]
再对上式作一次前向差分即得
\[
\mathcal L_q-2\mathcal L_{q-1}+\mathcal L_{q-2}
=
2\log F_q.
\]
证毕。
\end{proof}

\begin{theorem}[归一化主子式曲率的 Ramanujan--divisor 层析]\label{thm:conclusion-lucas-hankel-curvature-ramanujan-tomography}
\leanverified{paper\_conclusion\_lucas\_hankel\_curvature\_ramanujan\_tomography}
在推论 \ref{cor:conclusion-lucas-log-fibonacci-from-hankel-curvature} 的实分支假设下，对每个 \(r\ge 1\) 定义形式幂级数
\[
\mathfrak F_r(z):=\sum_{q\ge 1}c_r(q)\log F_q\,z^q,
\]
其中 \(c_r(\cdot)\) 为 Ramanujan 和。则有
\[
\mathfrak F_r(z)
=
\frac12\sum_{q\ge 2}c_r(q)\,\nabla^2\mathcal L_q\,z^q.
\]
并且对每个 \(r\ge 1\) 成立 divisor-sieve 恒等式
\[
\frac1r\sum_{d\mid r}\mathfrak F_d(z)
=
\sum_{\ell\ge 1}\log F_{r\ell}\,z^{r\ell}.
\]
\end{theorem}
\begin{proof}
第一式由推论 \ref{cor:conclusion-lucas-log-fibonacci-from-hankel-curvature} 的
\(
\nabla^2\mathcal L_q=2\log F_q
\)
逐项代入即得。对第二式，
\[
\frac1r\sum_{d\mid r}\mathfrak F_d(z)
=
\sum_{q\ge 1}
\left(\frac1r\sum_{d\mid r}c_d(q)\right)\log F_q\,z^q.
\]
经典恒等式
\[
\sum_{d\mid r}c_d(q)=r\,\mathbf 1_{\{r\mid q\}}
\]
遂给出
\[
\frac1r\sum_{d\mid r}\mathfrak F_d(z)
=
\sum_{\substack{q\ge 1\\ r\mid q}}\log F_q\,z^q
=
\sum_{\ell\ge 1}\log F_{r\ell}\,z^{r\ell}.
\]
证毕。
\end{proof}

\begin{theorem}[归一化 \texorpdfstring{$p$}{p}-进刚性的离散 Newton 势表达]\label{thm:conclusion-lucas-normalized-stiffness-newton-potential}
\leanverified{paper\_conclusion\_lucas\_normalized\_stiffness\_newton\_potential}
设 \(p\) 为任意素数，并假设 \(F_d\neq 0\) 对 \(1\le d\le q\) 成立。定义
\[
\widetilde{\mathrm{Stiff}}_p(q):=v_p(\mathcal H_q).
\]
则对一切 \(q\ge 1\) 有
\[
\widetilde{\mathrm{Stiff}}_p(q)
=
2\sum_{d=1}^{q}(q+1-d)\,v_p(F_d),
\]
并且对一切 \(q\ge 2\) 有精确二阶差分律
\[
\nabla^2\widetilde{\mathrm{Stiff}}_p(q)=2\,v_p(F_q).
\]
\end{theorem}
\begin{proof}
由定理 \ref{thm:conclusion-lucas-normalized-hankel-discrete-curvature}，
\[
\mathcal H_q=\prod_{d=1}^{q}F_d^{\,2(q+1-d)}.
\]
对上式取 \(p\)-进赋值即得
\[
v_p(\mathcal H_q)
=
2\sum_{d=1}^{q}(q+1-d)\,v_p(F_d).
\]
于是
\[
\nabla \widetilde{\mathrm{Stiff}}_p(q)
=
\widetilde{\mathrm{Stiff}}_p(q)-\widetilde{\mathrm{Stiff}}_p(q-1)
=
2\sum_{d=1}^{q}v_p(F_d),
\]
再作一次差分便得到
\[
\nabla^2\widetilde{\mathrm{Stiff}}_p(q)=2\,v_p(F_q).
\]
证毕。
\end{proof}

\begin{corollary}[归一化 \texorpdfstring{$p$}{p}-进刚性的首次出现原则]\label{cor:conclusion-lucas-stiffness-rank-of-apparition}
\leanverified{paper\_conclusion\_lucas\_stiffness\_rank\_of\_apparition}
若
\[
\rho_p:=\min\{n\ge 1:\ p\mid F_n\}
\]
存在，则对一切 \(q<\rho_p\) 有
\[
\widetilde{\mathrm{Stiff}}_p(q)=0,
\]
并且
\[
\widetilde{\mathrm{Stiff}}_p(\rho_p)=2\,v_p(F_{\rho_p})>0,
\qquad
\min\{q\ge 1:\ \nabla^2\widetilde{\mathrm{Stiff}}_p(q)\neq 0\}=\rho_p.
\]
\end{corollary}
\begin{proof}
由 \(\rho_p\) 的定义，当 \(q<\rho_p\) 时 \(v_p(F_d)=0\) 对所有 \(1\le d\le q\) 成立。定理 \ref{thm:conclusion-lucas-normalized-stiffness-newton-potential} 立即给出
\(
\widetilde{\mathrm{Stiff}}_p(q)=0
\)。
当 \(q=\rho_p\) 时，同一定理给出
\[
\widetilde{\mathrm{Stiff}}_p(\rho_p)
=
2\sum_{d=1}^{\rho_p}(\rho_p+1-d)\,v_p(F_d)
=
2\,v_p(F_{\rho_p})>0,
\]
因为此前各项均为零。最后，由
\(
\nabla^2\widetilde{\mathrm{Stiff}}_p(q)=2v_p(F_q)
\)
可知其首次非零位置恰为 \(\rho_p\)。证毕。
\end{proof}

\begin{proposition}[residue-lumping 的 Ramanujan 阴影公式]\label{prop:conclusion-charp-residue-lumping-ramanujan-shadows}
沿用定理 \ref{thm:conclusion-charp-hankel-rank-classification} 的记号。设
\[
\operatorname{char}(K)=p>q,\qquad
m:=\operatorname{ord}(s)\le q,
\]
并记
\[
C_a:=\sum_{\substack{0\le k\le q\\ k\equiv a\ (\mathrm{mod}\ m)}}\binom{q}{k},
\qquad
M:=\#\{a\in\{0,\dots,m-1\}:C_a\neq 0\}.
\]
对每个 \(r\mid m\)，令
\[
\mathcal B_r^{(q,m)}:=\sum_{a=0}^{m-1}c_r(a)\,C_a,
\qquad
\omega_r:=e^{2\pi i/r},
\qquad
U(r):=(\ZZ/r\ZZ)^\times.
\]
则有
\[
\mathcal B_r^{(q,m)}
=
\sum_{k=0}^{q}\binom{q}{k}c_r(k)
=
\sum_{j\in U(r)}(1+\omega_r^j)^q,
\]
并且成立 divisor-sieve 累积恢复公式
\[
\frac1r\sum_{d\mid r}\mathcal B_d^{(q,m)}
=
\sum_{\substack{0\le a<m\\ r\mid a}}C_a.
\]
\end{proposition}
\begin{proof}
因为 \(r\mid m\)，故 \(k\equiv a\pmod m\) 蕴含 \(k\equiv a\pmod r\)，从而 \(c_r(k)=c_r(a)\)。于是
\[
\mathcal B_r^{(q,m)}
=
\sum_{a=0}^{m-1}c_r(a)\sum_{\substack{0\le k\le q\\ k\equiv a\ (\mathrm{mod}\ m)}}\binom{q}{k}
=
\sum_{k=0}^{q}\binom{q}{k}c_r(k).
\]
再用
\[
c_r(k)=\sum_{j\in U(r)}\omega_r^{jk},
\]
可得
\[
\mathcal B_r^{(q,m)}
=
\sum_{j\in U(r)}\sum_{k=0}^{q}\binom{q}{k}\omega_r^{jk}
=
\sum_{j\in U(r)}(1+\omega_r^j)^q.
\]
另一方面，
\[
\frac1r\sum_{d\mid r}\mathcal B_d^{(q,m)}
=
\sum_{a=0}^{m-1}
\left(\frac1r\sum_{d\mid r}c_d(a)\right)C_a.
\]
再用
\[
\sum_{d\mid r}c_d(a)=r\,\mathbf 1_{\{r\mid a\}}
\]
即可得到
\[
\frac1r\sum_{d\mid r}\mathcal B_d^{(q,m)}
=
\sum_{\substack{0\le a<m\\ r\mid a}}C_a.
\]
证毕。
\end{proof}

\begin{conjecture}[Lucas--Hankel 刚性的最小 Ramanujan 图册]\label{conj:conclusion-lucas-hankel-minimal-ramanujan-atlas}
对定理 \ref{thm:conclusion-lucas-normalized-hankel-discrete-curvature} 与定理 \ref{thm:conclusion-charp-hankel-rank-classification} 所覆盖的 Lucas/Fibonacci 型输入，以下两族 Ramanujan 坐标应共同构成一套极小完备图册：
\begin{enumerate}
  \item Archimedean/\(p\)-adic 刚性图册
  \[
  \{\mathfrak F_r(z)\}_{r\ge 1}
  \qquad\text{或等价地}\qquad
  \{\widetilde{\mathrm{Stiff}}_p(q)\}_{p,\;q},
  \]
  它完全控制归一化 Hankel 主子式的离散曲率；
  \item 特征 \(p\) 塌缩图册
  \[
  \{\mathcal B_r^{(q,m)}\}_{r\mid m},
  \]
  它完全控制 residue-class 合并所诱导的秩缺陷。
\end{enumerate}
更具体地，对任意固定 \(q\) 与任意特征 \(p>q\)，定理 \ref{thm:conclusion-charp-hankel-rank-classification} 中的秩缺陷
\[
(q+1)-M
\]
（其中 \(M=\#\{a\in\{0,\dots,m-1\}:C_a\neq 0\}\)）应只依赖于
\[
\{\mathcal B_r^{(q,m)}\}_{r\mid m},
\]
而不依赖于更细的有序 \(m\)-元组 \((C_0,\dots,C_{m-1})\)。
\end{conjecture}

\endinput
